problem:  
Let \((S^{n+1},g)\) be a smooth Riemannian manifold whose Ricci curvature satisfies  
\[
\mathrm{Ric}_g \ge n(1-\varepsilon)g, \quad \text{with } \varepsilon>0 \text{ small,}
\]
and \(\mathrm{vol}_g(S^{n+1}) = \mathrm{vol}_{g_{\mathrm{round}}}(S^{n+1})\).  
Let \(M^n \subset (S^{n+1},g)\) be a closed, embedded, two-sided **min–max** minimal hypersurface of *lowest possible Morse index* among all nontrivial min–max hypersurfaces produced by the Almgren–Pitts construction for \(g\).  

**Research question:**  
For sufficiently small \(\varepsilon>0\), must \(M\) be diffeomorphic to the equatorial sphere?  
Moreover, does there exist a universal constant \(c_n>0\) such that the first **nonzero eigenvalue** \(\lambda_1(M)\) of the Laplace–Beltrami operator on \(M\) satisfies  
\[
|\lambda_1(M) - n| \le c_n\,\varepsilon^p,
\]
for some exponent \(p>0\) independent of \(M\)?  

In other words, can one detect the “near-roundness” of \((S^{n+1},g)\) purely from the spectral and index-theoretic properties of its *lowest-index min–max minimal hypersurface*, under only a Ricci lower bound and volume normalization?

---

Why is it a "good" problem:  
This problem corrects the logical inconsistency of the previous formulation and pushes the conceptual boundary beyond classical rigidity theorems. It discards area-minimizing assumptions (which would force stability) and instead anchors the discussion in the **min–max framework**, where index is the primary variational invariant known to survive weak geometric limits.  

The problem blends **three regimes of geometric analysis** that have rarely been integrated:
1. **Ricci curvature rigidity** (which controls large-scale geometry but not fine sectional behavior),  
2. **Spectral geometry of the induced Laplacian**, probing extrinsic curvature through intrinsic eigenvalues, and  
3. **Variational index theory for min–max hypersurfaces**, capturing global topological and analytic aspects of minimal surface theory.  

It extends the Urbano‑type index gap question to a broader, Ricci‑controlled but purely analytic setting: if the ambient sphere is almost Einstein, does the *first min–max minimal hypersurface* already “feel” the round geometry through its index and intrinsic spectrum?  

This formulation is simple but profound. It asks whether the deep rigidity phenomena known for exactly round metrics are *quantitatively stable* under Ricci bounds — a question whose resolution likely requires new analytic tools linking Ricci pinching, spectral concentration, and the Almgren–Pitts min–max hierarchy. The problem therefore exposes a fertile and unexplored interface between geometric measure theory, spectral analysis, and global curvature geometry.